\chapter{Weighted variation and stability of translators}\label{sec:translator-variation}
Let $X:\Sigma^n\to\mathbb R^{n+1}$ be a smooth immersed hypersurface, $n\ge1$. For two-sided statements choose a smooth unit normal $\nu$, set $A=D\nu$ and $H=\operatorname{tr}A$, and use $\vec H=-H\nu$. Fix a constant velocity $v$. The translator equation is $H+\langle v,\nu\rangle=0$. We derive the local variation formulas in \cite{HIMW19}, Proposition 3.2, pp.~8--9, converting its downward velocity to this general convention. Compact support is always intrinsic and away from any boundary; finite total weighted area is not assumed.

\begin{definition}[Translator weight and Jacobi operator]\label{def:translator-jacobi-weight}
Put $h=\langle X,v\rangle$ and $\rho_v=e^h$. Define
\[
 \mathcal D_v u=\Delta_\Sigma u+\langle v^T,\nabla u\rangle,
 \qquad J_vu=\mathcal D_vu+|A|^2u,
 \qquad Q_v(\phi)=\int_\Sigma(|\nabla\phi|^2-|A|^2\phi^2)\rho_v\,dA.
\]
Weighted stability means $Q_v(\phi)\ge0$ for every smooth compactly supported interior test. This is a second variation for fixed $v$; it is distinct from dynamical stability under mean curvature flow and from ordinary minimal-surface stability. Compare \cite{HIMW19}, equations (3.5)--(3.7).
\end{definition}

\begin{lemma}[Translator drift in divergence form]\label{lem:translator-weighted-parts}
On any smooth immersed hypersurface, whether translating or not,
\[
 \mathcal D_vu=\rho_v^{-1}\operatorname{div}_\Sigma(\rho_v\nabla u),
 \qquad \int\phi\mathcal D_vu\,\rho_v=-\int\langle\nabla\phi,\nabla u\rangle\rho_v
\]
for smooth $u$ and compactly supported smooth interior $\phi$.
\end{lemma}
\begin{proof}
The tangential gradient of $h$ is $v^T$, hence $\nabla\rho_v=\rho_vv^T$. Expand the divergence and integrate the compactly supported vector field $\phi\rho_v\nabla u$. This proves both identities without an assumption on the total weighted area.
\end{proof}

\begin{proposition}[Linearization and second variation for fixed velocity]\label{prop:translator-jacobi-linearization}
At a two-sided translator, the derivative of $H+\langle v,\nu\rangle$ under normal speed $f\nu$ is $-J_vf$. The second variation of $\int e^{\langle X,v\rangle}dA$ under compactly supported normal variations is $Q_v(f)$.
\end{proposition}
\begin{proof}
The scalar variation formula, Proposition \ref{prop:riemannian-mean-curvature-variation}, gives $H'=-(\Delta+|A|^2)f$ in Euclidean space. Differentiating the orthogonality and unit-length equations for the normal gives $\nu'=-\nabla f$. Thus $\langle v,\nu\rangle'=-\langle v^T,\nabla f\rangle$, proving the linearization.

The weighted first variation of Proposition \ref{prop:translator-weighted-area-example} is $\int\rho_v(H+\langle v,\nu\rangle)f$. At a translator, every term differentiating the weight, speed or area measure is multiplied by the vanishing Euler expression. The remaining term is $-\int\rho_vfJ_vf=Q_v(f)$ by Lemma \ref{lem:translator-weighted-parts}. This derives the sign of the formula rather than inferring it from the source's opposite velocity convention.
\end{proof}

\begin{proposition}[Translation angles and mean curvature are Jacobi fields]\label{prop:translator-angle-jacobi}
For every constant ambient vector $a$, the normal component $w_a=\langle a,\nu\rangle$ satisfies $J_vw_a=0$ on a two-sided translator. In particular $J_vH=0$.
\end{proposition}
\begin{proof}
Ambient translations preserve both the Euclidean curvature and the equation with fixed $v$. Their normal speed is $w_a$. The tangential part differentiates the identically zero Euler expression, so Proposition \ref{prop:translator-jacobi-linearization} gives $J_vw_a=0$. Since $H=-w_v$, the last assertion follows. This argument concerns the translator equation even though translating the surface may multiply its weighted area by a constant.
\end{proof}

\begin{proposition}[Height and exponential drift identities]\label{prop:translator-height-drift}
On a translator,
\[
 \mathcal D_vX=v,\qquad \mathcal D_vh=|v|^2,\qquad
 |\nabla h|^2=|v|^2-H^2,
 \qquad J_v(e^{-h})=(|A|^2-H^2)e^{-h}.
\]
The first two identities do not require two-sidedness; the latter two use the scalar notation on a two-sided hypersurface.
\end{proposition}
\begin{proof}
The immersion identity gives $\Delta X=\vec H=v^\perp$, and the drift term applied to $X$ is $v^T$. Pair their sum with $v$ to obtain the first two formulas. Since $v^\perp=-H\nu$, the orthogonal decomposition gives the gradient identity. The chain rule now yields
\[
 \mathcal D_ve^{-h}=e^{-h}(-\mathcal D_vh+|\nabla h|^2)=-H^2e^{-h}.
\]
Adding the potential proves the last formula. These identities are pointwise; integration over the whole surface would need additional integrability.
\end{proof}

\begin{proposition}[A weighted ground-state identity]\label{prop:translator-ground-state}
Let $w>0$ be smooth on a two-sided translator. For every smooth compactly supported interior $\phi$,
\[
 Q_v(\phi)=\int w^2\left|\nabla\left(\frac\phi w\right)\right|^2\rho_v
 -\int\frac{J_vw}{w}\phi^2\rho_v.
\]
Consequently a positive solution or supersolution $J_vw\le0$ proves weighted stability.
\end{proposition}
\begin{proof}
Set $\eta=\phi/w$. Expanding $|\nabla(w\eta)|^2$ leaves $w^2|\nabla\eta|^2$ and the terms
$\langle\nabla w,\nabla(w\eta^2)\rangle-|A|^2w^2\eta^2$. Integrate the latter by Lemma \ref{lem:translator-weighted-parts}; their sum is $-w\eta^2J_vw$. This is exactly the stated identity. Positivity of $w$ makes $\eta$ a legitimate compactly supported smooth test even when $w$ tends to zero at infinity.
\end{proof}

\begin{corollary}[Translating graphs are weighted stable]\label{cor:translating-graph-weighted-stability}
Every smooth translating graph over a nonempty connected open subset of $\mathbb R^n$ is weighted stable. Its quadratic form is strictly positive on every nonzero compactly supported smooth test, although its variational spectral bottom need not be positive.
\end{corollary}
\begin{proof}
The upward normal has positive vertical angle $w=(1+|Du|^2)^{-1/2}$, and Proposition \ref{prop:translator-angle-jacobi} gives $J_vw=0$. Apply Proposition \ref{prop:translator-ground-state}. If the quadratic form vanishes, $\phi/w$ has zero gradient, so connectedness makes it constant. The graph, as a manifold without its domain boundary, is noncompact; a nonzero constant multiple of $w>0$ cannot have compact support. Thus $\phi=0$. A sequence of compact tests can still have Rayleigh quotients tending to zero, as the grim-reaper example will show. No uniform bound for $w$ away from zero is needed.
\end{proof}

\begin{proposition}[A complete translator has infinite translator-weighted area]\label{prop:translator-infinite-weighted-area}
Let a nonempty connected complete boundaryless immersed hypersurface translate with $v\ne0$. Then $\int_\Sigma e^{\langle X,v\rangle}\,dA=\infty$, regardless of properness.
\end{proposition}
\begin{proof}
Assume the integral finite. Choose intrinsic Lipschitz cutoffs $0\le\eta_R\le1$, equal to one on the radius-$R$ intrinsic ball, supported in the radius-$2R$ ball, with $|\nabla\eta_R|\le2/R$. Completeness makes the supports compact. The integration-by-parts formula extends to these tests by smooth approximation on compact subsets. Proposition \ref{prop:translator-height-drift} and $|\nabla h|\le|v|$ give
\[
 |v|^2\int\eta_R\rho_v=-\int\langle\nabla\eta_R,\nabla h\rangle\rho_v,
 \qquad
 \left|\int\langle\nabla\eta_R,\nabla h\rangle\rho_v\right|
 \le\frac{2|v|}{R}\int\rho_v.
\]
Dominated convergence makes the left side tend to a strictly positive number, while the right side tends to zero. This contradiction proves the claim. All cutoffs are intrinsic, so no unmentioned properness assumption enters.
\end{proof}

\begin{proposition}[A weighted gradient gap and boundary-volume inequality]\label{prop:translator-weighted-gradient-gap}
On a translator with $v\ne0$, every compactly supported smooth interior $\phi$ satisfies
\[
 \int|\nabla\phi|^2\rho_v\ge\frac{|v|^2}{4}\int\phi^2\rho_v.
\]
For a relatively compact smooth domain $D$ in the hypersurface,
\[
 \int_{\partial D}\rho_v\ge |v|\int_D\rho_v.
\]
These are estimates for the drift-gradient energy; subtracting $|A|^2$ to form $Q_v$ can remove the positive gap.
\end{proposition}
\begin{proof}
Test $\mathcal D_vh=|v|^2$ against $\phi^2$ and integrate by parts:
\[
 |v|^2\int\phi^2\rho_v=-2\int\phi\langle\nabla\phi,\nabla h\rangle\rho_v
 \le2|v|\left(\int\phi^2\rho_v\right)^{1/2}
             \left(\int|\nabla\phi|^2\rho_v\right)^{1/2}.
\]
Division and squaring prove the first estimate, including the zero-test case separately. For the second, the divergence formula gives
$|v|^2\int_D\rho_v=\int_{\partial D}\rho_v\langle v^T,\eta\rangle\le|v|\int_{\partial D}\rho_v$.
The outward conormal and the pointwise bound $|v^T|\le|v|$ fix the sign and constant.
\end{proof}

\begin{lemma}[A positive Jacobi field controls the weighted square-integrable kernel]\label{lem:translator-l2-kernel-ground-state}
Let $\Sigma$ be connected, complete, boundaryless and two-sided, and suppose $w>0$ satisfies $J_vw=0$. Every smooth solution $J_vu=0$ with $\int u^2\rho_v<\infty$ is a constant multiple of $w$. Thus the weighted square-integrable kernel is zero if $w\notin L^2(\rho_vdA)$, and is spanned by $w$ if $w\in L^2(\rho_vdA)$.
\end{lemma}
\begin{proof}
For an intrinsic cutoff $\eta_R$ as in Proposition \ref{prop:translator-infinite-weighted-area}, integration by parts against $\eta_R^2u$ gives the exact identity
\[
 Q_v(\eta_Ru)=\int u^2|\nabla\eta_R|^2\rho_v.
\]
Smooth approximation justifies Lipschitz cutoffs. Combining with Proposition \ref{prop:translator-ground-state} yields
\[
 \int w^2\left|\nabla\left(\frac{\eta_Ru}{w}\right)\right|^2\rho_v
 \le\frac4{R^2}\int u^2\rho_v.
\]
On any fixed compact subset, $\eta_R=1$ for large $R$. Letting $R$ tend to infinity therefore shows $\nabla(u/w)=0$ there. Exhaustion and connectedness give $u=cw$. No global integrability of $|A|^2u^2$ was assumed: both cancellations occur with compactly supported tests before the limit.
\end{proof}
